\chapter{Einführung in diese Vorlage}
\pagenumbering{arabic}
\section{Tipps}
\begin{itemize}
 \item Auf Overleaf kann diese Vorlage als .zip komprimierter Ordner importiert werden.
 \item \qt{Anführungszeichen} können mit \verb"\qt{Text}"
 eingefügt werden.
 \item Bücher und andere Quellen werden nur in der Bibliographie angezeigt, wenn auch auf sie verwiesen wird.
 \item Der Satzspiegel der Seiten ist nicht zentriert, da sich links ein Binderand befindet. Dieser kann in der Präambel konfiguriert werden.
 \item Die Art des Absatzes zwischen Abschnitten kann in der Präambel geändert werden.
 \item Es empfiehlt sich alle Label nach Typ zu strukturieren. Alle Bilder können zum Beispiel mit \texttt{fig:} begonnen werden.
 \item Sollten Fragen entstehen, lohnt sich zuerst ein Blick in die Präambel.
 \item Die Abstände der Seitenränder folgen denen der Seitenränder gängiger Textverarbeitungsprogramme und sind einzuhalten.
 \item Es wird anderthalbfacher Zeilenabstand verwendet.
 \item Die Schriftgröße beträgt 12 pt.
 \item Blocksatz und Silbentrennung sind obligatorisch.
\end{itemize}

\section{Text}
Dies ist ein Absatz mit Text.

Hier beginnt ein neuer Absatz. Sollte \LaTeX~ein langes Wort nicht kennen, kann mit \verb;\-; ein Bindestrich eingefügt werden, der nur bei Bedarf verwendet wird. 
Der Gedankenstrich (--) wird durch doppelte Eingabe des Bindestrichs \verb;--; realisiert.\\
Eine neue Zeile beginnt hier. Ein Zeilenumbruch kann unterbunden werden, indem der Befehl \verb;\mbox{}; verwendet wird.

\section{Bilder}

\begin{figure}[ht]
 \centering
 \includegraphics[width=0.3\textwidth]{bilder/czg.pdf}
 \quelle{Bildquelle} %dieser Befehl muss direkt unter "includegraphics" stehen
 \caption{Das Schullogo des Carl-Zeiss-Gymnasiums in Jena}
 \label{fig:logo}
\end{figure}


Hier steht ein Text mit Verweis auf das Bild \ref{fig:logo}, welches sich mit nur einem Mausklick findet.

Figuren lassen sich mit dem \texttt{float} Paket durch die \texttt{[H]} Option fest auf der Seite platzieren.

Ein abweichender Eintrag zu einer Abbildung für das Abbildungsverzeichnis kann bei Bedarf in eckigen Klammern bei der Beschriftung angegeben werden.

In jedem Fall ist eine Quellenangabe für die Abbildungen vorzunehmen. 
Diese kann mit dem Befehl \verb;\quelle{}; direkt unter \verb;\includegraphics; eingefügt werden.

Jede Abbildung hat eine Abbildungsunterschrift, die sich aus einer Nummer und einer kurzen Erklärung zusammensetzt.

\section{Tabellen}
Tabellen lassen sich sehr einfach über die Seite \url{https://www.tablesgenerator.com/} erstellen. Jede Tabelle hat eine Überschrift und unterhalb gegebenenfalls einen Quellennachweis.
\begin{table}[h]
\centering
\caption{Dies ist eine Tabelle}
\begin{tabular}{|l|l|}
\hline
Spalte 1 & Spalte 2 \\ \hline\hline
Inhalt   & Inhalt   \\ \hline
Inhalt   & Inhalt   \\ \hline
\end{tabular}\\\vspace{.1cm}
{Quelle: Quelle der Tabelle}
\end{table}


\section{Bücher und Internetquellen}
Auf Bücher und Internetquellen lässt sich sehr leicht in Fußnoten referenzieren.\\ 
Viel über grundlegende Programmierung lernt man in dem Buch von Donald E. Knuth.\footnote{Vgl. \cite{book:knuth97}.} \qt{Dieses Zitat benötigt eine Seitenangabe.}\footnote{S. \cite[S. 57]{book:knuth97}.} 
Eine beliebte Internetseite sind die Lexika des Spektrum-Verlags\footnote{Vgl. \cite{online:spektrumlex}.}, die für alle MINT-Fächer wissenschaftliche Nachschlagewerke bereitstellen.
Es wird in der Fußnote automatisch auf das Literatur- beziehungsweise das Internetquellenverzeichnis verwiesen. 
Sollte eine Darstellung der Quellen nach ISO 690 bevorzugt werden, kann diese in der Präambel ausgewählt werden.

Im Literaturverzeichnis am Ende der Arbeit werden alle genutzten Werke dann alphabetisch sortiert und genannt.

Nutzt für eine Orientierung bei der Angabe von Literatur die Broschüre auf der Homepage und einigt euch in der Gruppe auf ein gemeinsames Vorgehen.

\section{Quelltext}
Dies: \pythoninline{class GuteKlasse} ist Quelltext im Text und hier \ref{listing:code} in einer Figur als Listing.\\
Quelltext aus einem Source-File einzubinden ist mit
\begin{verbatim}
 \pythonexternal{file.py}
\end{verbatim}
auch möglich.

\begin{figure}[h!]
\centering
\begin{python}
class MyClass(SuperClass):
    def __init__(self):
        print("Hello World")
\end{python}
\caption{Quellcodelisting in einer Figur}
\label{listing:code}
\end{figure}